% part: intuitionistic-logic
% chapter: propositions-as-types
% section: introduction

\documentclass[../../../include/open-logic-section]{subfiles}

\begin{document}

\olfileid{int}{pty}{int}

\olsection{مقدمة}

حساب لامبدا نموذج بسيط للحوسبة. وهو يعمل على حدود مبنية من المتغيرات. فإذا
كان~$N$ و$M$ حدين، كان~$(NM)$ حدًا أيضًا (``تطبيق~$N$ على~$M$''). وإذا كان
$N$ حدًا، كان~$\lambd[x][N]$ حدًا. وإذا احتوى~$N$ المتغير~$x$، فإن وقائع
$x$ المقابلة في~$\lambd[x][N]$ تكون مقيّدة بالمعامل~$\lambd[x]$، فلا تعود
حرة. ويقال إن حدًا من الصورة~$(\lambd[x][N]M)$ \emph{ينقبض بيتا} إلى
$\Subst{N}{M}{x}$، أي إلى نتيجة استبدال~$M$ بكل وقوع حر لـ~$x$ في~$N$.
ويتحول~$N$ بتحويل بيتا إلى~$N'$ إذا نتج~$N'$ من~$N$ بانقباض بيتا لحد جزئي،
ونكتب~$N\redone N'$. والحوسبة في حساب~$\lambd$ متتالية
$N\redone N_1\redone N_2\redone\dots$. فإذا انتهت المتتالية بحد~$N_k$ لا
يمكن انقباض بيتا لأي حد جزئي منه، سُمّي \emph{سويًا} أو صورة سوية لـ~$N$.
وفكر في الحوسبة في حساب~$\lambd$ بوصفها متتاليات كهذه. تتوقف الحوسبة إذا
نتجت صورة سوية، وتكون تلك الصورة ناتج الحوسبة. ويتبين أن هذا النموذج البسيط
تام تورنغ: فكل دالة جزئية عودية يمكن حسابها بحد في حساب~$\lambd$.

اكتشف هاسكل كاري ووليام هوارد، كل على حدة، تشابهًا وثيقًا بين الاستنباط
الطبيعي وحساب~$\lambd$. فالحد~$\lambd[x][N]$ دالة بحسب الحدس: المتغير~$x$
هو المعطى، و$N$ يبين كيف نحسب القيمة انطلاقًا منه. ويكون
$(\lambd[x][N]M)$ تطبيق الدالة~$\lambd[x][N]$ على المعطى~$M$، ويبين
انقباض بيتا كيفية تقييمه: نستبدل~$M$ بـ~$x$ في~$N$، ثم نواصل تقييم الحد
الناتج. فكر الآن في هذه الدوال بوصفها ذات مجال ومدى ثابتين. فإذا كان مجال
$\lambd[x][N]$ هو~$A$ ومداه~$B$، وجب أن نعد~$x$ متغيرًا لا يأخذ إلا
معطيات من~$A$، وأن نعد~$N$ منتجًا قيمًا من~$B$. ولذلك لا يكون
$(\lambd[x][N]M)$ ذا معنى إلا إذا كان~$\lambd[x][N]$ دالة~$A\to B$ وكان
$M\in A$. وإذا أضفنا هذه المعلومات إلى حساب~$\lambd$ حصلنا على حساب
$\lambd$ \emph{المنمّط}. وليست كل عبارة~$(\lambd[x][N]M)$ مشروعة في
الحساب المنمّط، بل تقتصر المشروعة على العبارات التي تتلاءم فيها ``أنماط''
$\lambd[x][N]$ و$M$؛ أي يكون نمط~$\lambd[x][N]$ هو~$A\to B$ ونمط~$M$
هو~$A$. ولا يكون نمط~$\lambd[x][N]$ هو~$A\to B$ إلا إذا كان نمط~$x$ هو
$A$ ونمط~$N$ هو~$B$. وبوجه عام ليست كل عبارة~$(M_1M_2)$ مشروعة، بل
تقتصر الحدود المشروعة منها على ما كان فيه نمط~$M_1$ هو~$A\to B$ ونمط
$M_2$ هو~$A$؛ وعندئذ يكون نمط~$(M_1M_2)$ هو~$B$.

تقابل قواعد التنميط هذه بوضوح~!!{derivation}s في الاستنباط الطبيعي، حيث
تمثّل الأنماط~!!{formula}s، وتقابل~!!{derivation}s نفسها حدود~$\lambd$.
فمثلًا، يقابل~!!a{derivation} لـ~$!A\lif !B$ حدًا
$\lambd[\typeof{x}{!A}][N]$ نمطه الدالي~$!A\lif !B$. وتقابل قواعد
الاستدلال في الاستنباط الطبيعي قواعد التنميط في حساب لامبدا المنمّط، مثلًا:
\begin{prooftree}
  \AxiomC{$\Discharge{!A}{x}$}
  \DeduceC{$!B$}
  \DischargeRule{\Intro\lif}{x}
  \UnaryInfC{$!A \lif !B$}
  \DisplayProof\bottomAlignProof
  \quad\text{تقابل}\quad
  \Axiom$x: !A \fCenter N: !B$
  \UnaryInf$ \fCenter \lambd[\typeof{x}{!A}][N] : !A \lif !B$
\end{prooftree}
وتعني القاعدة على اليمين أنه إذا كان نمط~$x$ هو~$!A$ ونمط~$N$ هو~$!B$،
كان نمط~$\lambd[\typeof{x}{!A}][N]$ هو~$!A\lif !B$.

تقابل قاعدة~$\Elim\lif$ قاعدة التنميط لحدود التطبيق، أي:
\[
  \AxiomC{$!A \lif !B$}
  \AxiomC{$!A$}
  \RightLabel{\Elim\lif}
  \BinaryInfC{$!B$}
  \DisplayProof
  \quad\text{تقابل}\quad
  \Axiom$\fCenter M_1: !A \lif !B$
  \Axiom$\fCenter M_2: !A$
  \BinaryInf$ \fCenter (M_1M_2) : !B$
  \DisplayProof
\]
إذا تبعت قاعدة~$\Intro\lif{}$ فورًا قاعدة~$\Elim\lif{}$، أمكن تبسيط
!!{derivation}:
\begin{gather*}
  \AxiomC{$\Discharge{!A}{x}$}
  \DeduceC{$!B$}
  \DischargeRule{\Intro\lif}{x}
  \UnaryInfC{$!A \lif !B$}
  \AxiomC{}
  \DeduceC{$!A$}
  \RightLabel{\Elim\lif}
  \BinaryInfC{$!B$}
  \DisplayProof
  \quad\redone\quad
  \AxiomC{}
  \DeduceC{$!A$}
  \DeduceC{$!B$}
  \DisplayProof
\end{gather*}
وهذا يقابل انقباض بيتا لحدود لامبدا:
\[
(\lambd[\typeof{x}{!A}][N]M) \quad\redone\quad
\Subst{N}{M}{x}.
\]

وتوجد مقابلات مماثلة بين قواعد~$\land$ والأنماط ``الضربية''، وبين قواعد
$\lor$ والأنماط ``الجمعية''.

تسمى هذه المقابلة بين حدود حساب لامبدا البسيط المنمّط و!!{derivation}s
الاستنباط الطبيعي مقابلة ``كاري--هوارد'' أو ``البراهين بوصفها برامج'' أو
``القضايا بوصفها أنماطًا''. وإلى جانب مقابلة~!!{formula}s (القضايا) للأنماط
والبراهين للحدود، يمكن تلخيص المقابلات هكذا:
\begin{center}
  \begin{tabular}{c c}
    المنطق & البرنامج \\
    \hline
    القضية/!!{formula} & النمط \\
    البرهان & الحد \\
    الفرض & المتغير \\
    الفرض المفرّغ & المتغير المقيّد \\
    الفرض غير المفرّغ & المتغير الحر \\
    $!A \lif !B$ & النمط الدالي \\
    $!A \land !B$ & النمط الضربي \\
    $!A \lor !B$ & النمط الجمعي \\
    $\lfalse$ & النمط الخالي \\
    \hline
  \end{tabular}
\end{center}

مقابلة كاري--هوارد إحدى ركائز مساعدات البرهان الآلية وفاحصات أنماط البرامج،
لأن فحص برهان يشهد لقضية، كما فعلنا أعلاه، لا يعدو فحص ما إذا كان برنامج
(حد) يحمل النمط المعلن.
\end{document}
